\documentclass[11pt,a4paper]{article}
\usepackage[T1]{fontenc}
\usepackage[utf8]{inputenc}
\usepackage{amsmath,amssymb,amsthm}
\usepackage[margin=1in]{geometry}
\usepackage[colorlinks=true,linkcolor=blue,urlcolor=blue]{hyperref}
\theoremstyle{plain}
\newtheorem{theorem}{Theorem}
\newtheorem{lemma}[theorem]{Lemma}
\newtheorem{proposition}[theorem]{Proposition}
\newtheorem{corollary}[theorem]{Corollary}
\theoremstyle{definition}
\newtheorem{remark}[theorem]{Remark}

\title{The empirical recurrence for OEIS A198703, proved by a two-line transfer matrix}
\author{Adrian Perez Fontelles\\ \small Independent researcher}
\date{31 August 2026}
\begin{document}
\maketitle

\begin{abstract}
OEIS A198703 counts arrays in which every cell is constrained by its neighbours on \emph{both}
sides, and which are moreover written in canonical form (``values introduced in row major
order''). Neither feature is visible to a one-line transfer matrix: the cell condition
reaches the row above and the row below, and the canonical-form clause is not local at
all. Both are dealt with. What is counted is equality patterns, so
$5!\,a(n)=\sum_i\binom{5}{i}D_{5-i}L_i(n)$ with $D_m$ the derangement numbers and
$L_i$ the count over an $i$-letter alphabet; and each $L_i$ is a walk count whose vertices
are ordered pairs of consecutive rows. Relabelling-invariance makes the chain strongly
lumpable, cutting $20515$ pairs down to $S=544$ states. The entry's empirical recurrence of
order $39$, contributed by R.~H.~Hardin, then follows from a finite exact computation.
\end{abstract}

\noindent\small 2020 Mathematics Subject Classification. 05A15, 05A18, 15B36.\normalsize

\section{The sequence and the conjecture}

OEIS A198703 is ``Number of nX3 0..4 arrays with values 0..4 introduced in row major order and each element equal to at least one horizontal or vertical neighbor.''. It has offset $1$ and begins
\[
1, 14, 185, 3893, 95209, 2654693, 77622801, 2332544112,\ \dots
\]
The entry carries the following comment:
\begin{quote}\small
Empirical: a(n) = 36*a(n-1) -a(n-2) -5355*a(n-3) +4939*a(n-4) +171547*a(n-5) +187210*a(n-6) -2080273*a(n-7) -6373170*a(n-8) +4618206*a(n-9) +52436433*a(n-10) +65874012*a(n-11) -142129836*a(n-12) -448353529*a(n-13) -92258782*a(n-14) +1182054840*a(n-15) +1645176287*a(n-16) -1092230888*a(n-17) -4897845688*a(n-18) -3393264569*a(n-19) +5453181946*a(n-20) +10604051822*a(n-21) +3233851655*a(n-22) -11263670956*a(n-23) -13806085667*a(n-24) +2193563289*a(n-25) +13819377251*a(n-26) +6644812359*a(n-27) -4919877474*a(n-28) -9065267335*a(n-29) +842958990*a(n-30) +3103300670*a(n-31) +1235389140*a(n-32) -446260304*a(n-33) -841232800*a(n-34) +553687700*a(n-35) -274224576*a(n-36) +45862920*a(n-37) +8532304*a(n-38) +2067648*a(n-39)
\end{quote}
As of the ``Last modified'' line on the live entry (Oct 03 2025 07:59:41, revision 9) the statement is
still recorded as empirical, and nothing on the entry records it as settled either way.

Written out, the claim is
\begin{equation}\label{eq:conj}
a(n)\;=\;36\,a(n-1) - a(n-2) - 5355\,a(n-3) + 4939\,a(n-4) + 171547\,a(n-5) + 187210\,a(n-6) - 2080273\,a(n-7) - 6373170\,a(n-8) + 4618206\,a(n-9) + 52436433\,a(n-10) + 65874012\,a(n-11) - 142129836\,a(n-12) - 448353529\,a(n-13) - 92258782\,a(n-14) + 1182054840\,a(n-15) + 1645176287\,a(n-16) - 1092230888\,a(n-17) - 4897845688\,a(n-18) - 3393264569\,a(n-19) + 5453181946\,a(n-20) + 10604051822\,a(n-21) + 3233851655\,a(n-22) - 11263670956\,a(n-23) - 13806085667\,a(n-24) + 2193563289\,a(n-25) + 13819377251\,a(n-26) + 6644812359\,a(n-27) - 4919877474\,a(n-28) - 9065267335\,a(n-29) + 842958990\,a(n-30) + 3103300670\,a(n-31) + 1235389140\,a(n-32) - 446260304\,a(n-33) - 841232800\,a(n-34) + 553687700\,a(n-35) - 274224576\,a(n-36) + 45862920\,a(n-37) + 8532304\,a(n-38) + 2067648\,a(n-39)
\end{equation}
for all $n>40$.

\section{The condition, and what the canonical-form clause counts}

The arrays have $3$ columns and $L=n$ rows. Write $x_{t,u}$ for the entry in
row $t$ and column $u$. with the neighbour offsets
\[
\mathcal N=\{(-1,0),\ (1,0),\ (0,-1),\ (0,1)\}
\]
written as (row shift, column shift), the requirement is
\begin{equation}\label{eq:loc}
\#\bigl\{(d_1,d_2)\in\mathcal N:\ (t+d_1,u+d_2)\ \text{lies in the array and its entry is }
equal to x_{t,u}\bigr\}\ \ge1
\end{equation}
for every cell $(t,u)$.

The condition is stated purely in equalities between entries, so it is unchanged by any
relabelling of the values: it is a property of the \emph{equality pattern} --- the partition
of the cells into classes of equal entries --- and not of the values themselves. The clause
``values introduced in row major order'' selects exactly one array from each relabelling
class, so with $K=5$ available values the entry counts admissible patterns with at most $K$
classes.

\begin{lemma}\label{lem:inv}
Let $L_i(n)$ count the arrays over an alphabet of $i$ letters that satisfy the cell
condition, with no canonical-form clause, and $N_j(n)$ the admissible patterns with exactly
$j$ classes. Then $L_i=\sum_j N_j\,i(i-1)\cdots(i-j+1)$ and
\[
5!\,\sum_{j\le 5}N_j(n)\;=\;\sum_{i=0}^{5}\binom{5}{i}D_{5-i}\,L_i(n),
\qquad D_m=m!\sum_{t=0}^{m}\frac{(-1)^t}{t!} .
\]
\end{lemma}

\begin{proof}
A labelled array satisfying the condition is a pattern together with an injection of its
classes into the alphabet, and a pattern with $j$ classes admits $i(i-1)\cdots(i-j+1)$ of
them; that is the first identity. Writing $i(i-1)\cdots(i-j+1)=j!\binom{i}{j}$ it reads
$L_i=\sum_j(j!N_j)\binom{i}{j}$, so binomial inversion gives
$j!N_j=\sum_i(-1)^{j-i}\binom{j}{i}L_i$. Summing over $j\le5$ and exchanging the
order of summation, the coefficient of $L_i$ is
$\sum_{t\le5-i}(-1)^t/(i!\,t!)$, which multiplied by $5!$ is
$\binom{5}{i}(5-i)!\sum_{t\le5-i}(-1)^t/t!=\binom{5}{i}D_{5-i}$.
\end{proof}

\section{Each $L_i$ is a walk count on pairs of rows}

Fix $i$ and let $V_i=\{0,\dots,i-1\}^{3}$. Because \eqref{eq:loc} at a cell of row
$t$ involves rows $t-1$, $t$ and $t+1$, a single row is not enough state. Take the
ordered pairs $(p,c)\in V_i\times V_i$ as vertices and declare $(p,c)\to(c,x)$ an edge when
every cell of the middle row $c$ satisfies \eqref{eq:loc} with $p$ above it and $x$
below it. Let $P_i$ be the adjacency matrix, $\mathbf{s}_i$ the indicator of the pairs
$(p,c)$ whose first row $p$ satisfies the condition with no row above, and
$\mathbf{e}_i$ the indicator of the pairs whose second row $c$ satisfies it with no
row below.

\begin{lemma}\label{lem:walk}
For $L\ge2$, the arrays of $L$ rows satisfying the condition are exactly the walks
$(r^{(1)},r^{(2)})\to(r^{(2)},r^{(3)})\to\cdots\to(r^{(L-1)},r^{(L)})$ of length
$L-2$ that start at a vertex marked by $\mathbf{s}_i$ and end at one marked by
$\mathbf{e}_i$, so
$L_i=\mathbf{s}_i^{\!\top}P_i^{\,L-2}\mathbf{e}_i$.
\end{lemma}

\begin{proof}
A walk of that form is exactly a sequence of rows $r^{(1)},\dots,r^{(L)}$, since
consecutive vertices overlap in one row. In such a walk the condition at the cells of
row $t$ is tested once, at the step leaving $(r^{(t-1)},r^{(t)})$, and with the
correct neighbours $r^{(t-1)}$ and $r^{(t+1)}$; the cells of row $1$ are tested by
$\mathbf{s}_i$, which uses the condition with the row above absent, and those of row
$L$ by $\mathbf{e}_i$. So every cell is tested exactly once and with its true
neighbourhood, which is the claim.
\end{proof}

\begin{lemma}\label{lem:lump}
Write $\pi(p,c)$ for the equality pattern of the pair. If $\pi(p,c)=\pi(p',c')$ then for
every pattern $Q$ the two vertices have the same number of successors of pattern $Q$, and
$\mathbf{s}_i$, $\mathbf{e}_i$ take the same value at both. Hence the chain is strongly
lumpable over $\pi$ and $L_i$ is computed on the quotient, whose states are the set
partitions of the $2\cdot3$ cells of a pair into at most $i$ classes.
\end{lemma}

\begin{proof}
Equal patterns means $(p',c')=(\sigma(p),\sigma(c))$ for a permutation $\sigma$ of the
alphabet. The edge condition and the two boundary conditions are statements about equalities
between entries, so they are invariant under $\sigma$, and $x\mapsto\sigma(x)$ is a bijection
between the successor sets preserving patterns. The number of pairs with a given pattern $P$
is $e_P=i(i-1)\cdots(i-|P|+1)$, the number of injections of its classes into the alphabet,
which supplies the weights on the quotient.
\end{proof}

Assembling the $5$ lumped chains block-diagonally into $N$ of size $S=544$, with the
weights of Lemma~\ref{lem:inv} folded into the start vector $\mathbf{v}$ and the
end-of-array indicator into $\mathbf{u}$,
\begin{equation}\label{eq:walk2}
5!\;a(n)\;=\;\mathbf{v}^{\!\top}N^{\,L-2}\mathbf{u},\qquad L=n .
\end{equation}

\section{The criterion, and the computation}

Let $q(t)=t^{39}-\sum_i c_i\,t^{39-i}$ be the characteristic polynomial of
\eqref{eq:conj} and put $G=N^{1}$, so that consecutive values of $n$ advance the
walk by $1$ rows.

\begin{lemma}\label{lem:crit}
With $h_j=G^{j}N^{o}\mathbf{u}$ for the offset $o$ that aligns the first admissible
index,
\[
5!\Bigl(a(n)-\sum_i c_i\,a(n-i)\Bigr)\;=\;\mathbf{v}^{\!\top}G^{\,j-39}\,
\Bigl(h_{39}-\sum_i c_i h_{39-i}\Bigr)
\]
for the corresponding $j$. Hence \eqref{eq:conj} holds from that point on if and only if
$\mathbf{v}^{\!\top}G^{m}w=0$ for every $m\ge0$, where
$w=h_{39}-\sum_i c_ih_{39-i}$; and it is enough to check $m=0,1,\dots,S-1$.
\end{lemma}

\begin{proof}
Substitute \eqref{eq:walk2} and factor. The scalar $5!$ is nonzero, so it does not
affect vanishing. For the last clause, Cayley--Hamilton applied to $G$ makes $G^{S}$ an
integer combination of $I,G,\dots,G^{S-1}$, so each $\mathbf{v}^{\!\top}G^{m}w$ with
$m\ge S$ is the same combination of the earlier values; if those vanish, so do all the rest.
\end{proof}

\begin{theorem}
The recurrence \eqref{eq:conj} holds for every $n>40$, which is the statement of
Section~1.
\end{theorem}

\begin{proof}
The lumped transition counts were built by taking one representative pair for each pattern,
enumerating every row $x$ over the alphabet, testing \eqref{eq:loc} on the middle
row, and sorting the results by the pattern of the new pair. The vector $w$ was formed by
$39$ applications of $G$ in exact integer arithmetic and $\mathbf{v}^{\!\top}G^{m}w$
evaluated for $m=0,1,\dots$, again exactly. Every one of those integers vanishes from the
index corresponding to $n=40+1$ onwards, and the run was continued until $S+1$
consecutive zeros had been seen, which by Lemma~\ref{lem:crit} settles every larger $m$.
\end{proof}

\section{Verification}

Three checks, each independent of the algebra above.

First, the whole construction was compared against the entry's DATA: the right-hand side of
\eqref{eq:walk2}, divided by $5!$, was evaluated at every one of the $16$
published indices and agrees exactly, the quotient coming out an integer each time. This
tests the modelling and not only the arithmetic --- the inversion of Lemma~\ref{lem:inv},
the reading of \eqref{eq:loc}, the treatment of the first and last rows, and the
alignment of the entry's index with $L=n$ would each show up as a mismatch.

Second, the lumped chain was checked against the unlumped one on the small shapes of this
family: the walk counts over the full alphabet agree with $\sum_P e_Pf_m(P)$ term by term.

Third, the recurrence \eqref{eq:conj} was evaluated on the published terms in exact integer
arithmetic with no matrices involved, and the annihilation test of Lemma~\ref{lem:crit} was
carried to the full Cayley--Hamilton bound rather than to a sampled prefix.

\begin{thebibliography}{9}
\bibitem{oeis} The OEIS Foundation, \emph{The On-Line Encyclopedia of Integer
Sequences}, \url{https://oeis.org}, 2026. Sequence A198703.
\bibitem{stanley} R.~P.~Stanley, \emph{Enumerative Combinatorics}, Volume 1, 2nd ed.,
Cambridge University Press, 2011. (Transfer-matrix method, Section 4.7; falling-factorial
inversion, Section 1.9.)
\bibitem{kemeny} J.~G.~Kemeny and J.~L.~Snell, \emph{Finite Markov Chains}, Springer,
1976. (Strong lumpability, Section 6.3.)
\bibitem{fs} P.~Flajolet and R.~Sedgewick, \emph{Analytic Combinatorics}, Cambridge
University Press, 2009.
\end{thebibliography}

\end{document}
